% Part: first-order-logic
% Chapter: tableaux
% Section: quantifier-rules

\documentclass[../../../include/open-logic-section]{subfiles}

\begin{document}

\olfileid{fol}{tab}{qrl}

\olsection{Kwantorregels}

\subsection{Regels voor $\lforall$}

\begin{defish}
\AxiomC{\sFmla{\True}{\lforall[x][!A(x)]}}
\RightLabel{\TRule{\True}{\forall}}
\UnaryInfC{\sFmla{\True}{!A(t)}}
\DisplayProof
\hfill
\AxiomC{\sFmla{\False}{\lforall[x][!A(x)]}}
\RightLabel{\TRule{\False}{\lforall}}
\UnaryInfC{\sFmla{\False}{!A(a)}}
\DisplayProof
\end{defish}

In \TRule{\True}{\lforall} is $t$ een gesloten term (dat wil zeggen:
een term zonder variabelen). In \TRule{\False}{\lforall} is
$a$~!!a{constant} die nergens mag voorkomen in de tak boven de regel
\TRule{\False}{\lforall}. We noemen $a$ de \emph{eigenvariabele} van
de inferentie \TRule{\False}{\forall}.\footnote{We gebruiken de term
``eigenvariabele'', hoewel $a$ in de bovenstaande regel !!a{constant}
is. Dit heeft historische redenen.}

\subsection{Regels voor $\lexists$}

\begin{defish}
\AxiomC{\sFmla{\True}{\lexists[x][!A(x)]}}
\RightLabel{\TRule{\True}{\lexists}}
\UnaryInfC{\sFmla{\True}{!A(a)}}
\DisplayProof
\hfill
\AxiomC{\sFmla{\False}{\lexists[x][!A(x)]}}
\RightLabel{\TRule{\False}{\lexists}}
\UnaryInfC{\sFmla{\False}{!A(t)}}
\DisplayProof
\end{defish}

Ook hier is $t$~een gesloten term, en is $a$~!!a{constant} die niet
voorkomt in de tak boven de regel~\TRule{\True}{\lexists}. We noemen
$a$ de \emph{eigenvariabele} van de inferentie
\TRule{\True}{\lexists}.

De voorwaarde dat een eigenvariabele niet voorkomt in de tak boven de
inferentie \TRule{\False}{\lforall} of \TRule{\True}{\lexists}, heet
de \emph{eigenvariabelevoorwaarde}.

\begin{explain}
Herinner u de afspraak dat we, wanneer $!A$ !!a{formula} is waarin de
!!{variable}~$x$ vrij voorkomt, dit aangeven door~$!A(x)$ te schrijven.
In dezelfde context is $!A(t)$ dan een afkorting voor
$\Subst{!A}{t}{x}$. We zouden de regel
\TRule{\False}{\lexists} dus ook als volgt kunnen schrijven:
\begin{prooftree}
  \AxiomC{\sFmla{\False}{\lexists[x][!A]}}
  \RightLabel{\TRule{\False}{\lexists}}
  \UnaryInfC{\sFmla{\False}{\Subst{!A}{t}{x}}}
\end{prooftree}
Merk op dat $t$ al in~$!A$ kan voorkomen; $!A$ kan bijvoorbeeld
$\Atom{\Obj P}{t,x}$ zijn. Het afleiden van
$\sFmla{\False}{\Atom{\Obj P}{t,t}}$ uit
$ \sFmla{\False}{\lexists[x][\Atom{\Obj P}{t,x}]}$ is dus een
correcte toepassing van~\TRule{\False}{\lexists}. De
eigenvariabelevoorwaarden in \TRule{\False}{\lforall}
en~\TRule{\True}{\lexists} vereisen echter dat de !!{constant}~$a$
niet in~$!A$ voorkomt. Men kan dus niet met een correcte toepassing van
de formule $\sFmla{\False}{\Atom{\Obj P}{a,a}}$ afleiden uit
$\sFmla{\False}{\lforall[x][\Atom{\Obj P}{a,x}]}$ met
$\TRule{\False}{\lforall}$.
\end{explain}

\begin{explain}
In \TRule{\True}{\lforall} en \TRule{\False}{\lexists} gelden geen
beperkingen voor de term~$t$. Bij de regels
\TRule{\True}{\lexists} en \TRule{\False}{\lforall} vereist de
eigenvariabelevoorwaarde daarentegen dat de !!{constant}~$a$ nergens
voorkomt in de takken boven de desbetreffende inferentie. Deze
voorwaarde is nodig om de correctheid van het systeem te waarborgen.
Zonder deze voorwaarde zou het volgende een gesloten !!{tableau} zijn
voor
$\lexists[x][\formula{A}(x)] \lif \lforall[x][\formula{A}(x)]$:
\begin{center}
\begin{tableau}{}
  [\sFmla{\False}{\lexists[x][\formula{A}(x)] \lif \lforall[x][\formula{A}(x)]}, just=\TAss
    [\sFmla{\True}{\lexists[x][\formula{A}(x)]},
      just={\TRule{\False}{\lif}[1]}
      [\sFmla{\False}{\lforall[x][\formula{A}(x)]},
        just={\TRule{\False}{\lif}[1]}
        [\sFmla{\True}{\formula{A}(a)},
          just={\TRule{\True}{\lexists}[2]}
          [\sFmla{\False}{\formula{A}(a)},
            just={\TRule{\False}{\lforall}[3]}, close]
        ]
      ]
    ]
  ]
\end{tableau}
\end{center}
$\lexists[x][\formula{A}(x)] \lif
\lforall[x][\formula{A}(x)]$ is echter niet geldig.
\end{explain}

\end{document}
